\part{Développements en langage C}
\label{part:dev_C}

\textbf{\textit{note : le codage en langage C est abandonné au profit du codage
en langage Python.}}

\section{Règles de codages}
   AIDA est pour l'instant un petit projet, qui ne concerne pas un
grand nombre de programmeurs. Nous n'avons donc pas mis en place
de règles de codages strictes. Ce n'est pas une raison pour écrire
du code incompréhensible, non-portable, ... Il est donc fortement
recommandé de connaître une des règles de codages existantes
(par exemple, le «GNU Coding standard» 
\footnote{\url{http://www.gnu.org/prep/standards_toc.html}})
, et de la suivre, sinon \`a la lettre, du 
moins dans les grands principes.

  Pour uniformiser la présentation et faciliter la (re)lecture
des sources, l'usage s'est répandu (et même systématisé à TRAMM)
d'automatiser la mise en forme des sources avec l'utilitaire \texttt{indent}.

  Les commentaires seront rédigés de manière à faciliter la
création d'une documentation technique par des outils automatiques.
Actuellement, la recommandation est d'utiliser l'outil
Doxygen\footnote{\url{https://www.doxygen.nl/}},
avec style de commentaires se
rapprochant de javadoc\footnote{\url{http://fr.wikipedia.org/wiki/Javadoc}}.
Historiquement, la documentation automatique des premiers programmes (la 
bibliothèque \texttt{rapt}, par exemple) était faite avec l'outil
\texttt{cextract/cextdoc}.
Certains commentaires «ancien style» persistent donc encore
dans les sources.


\section{La bibliothèque \texttt{rapt}}
\label{sec:rapt}
   Un certain nombre de fonctions utilisées fréquemment dans
l'environnement AIDA (lecture/écriture d'un vecteur, codage
ou décodage d'un "AcqId", ...) ont été rassemblées dans une 
bibliothèque, baptisée \textit{rapt}\footnote{Historiquement,
cela signifiait \textit{Research Aircraft Processing Tools}}.
Voici une présentation générale des différents groupes de fonctions
(une présentation détaillée de chaque fonction, avec le type
de chaque argument, et un court texte descriptif se trouve
dans le fichier \texttt{rapt.doc.ps} du répertoire 
\path{$AIDA_HOME/src/lib/rapt}).

\subsection{Fonctionnement général et traitement des erreurs}
  Cette bibliothèque fonctionne selon un système semblable à 
celui de la bibliothèque standard pour les appels systèmes en C. En cas
d'erreur, la fonction retourne une valeur invalide (indice
négatif, pointeur NULL, ...) et une variable globale 
précise le type d'erreur.

  Dans cette bibliothèque, la variable globale s'appelle \texttt{errapt},
et les fonctions \texttt{strerrapt()} et \texttt{perrapt()} permettent
d'obtenir, respectivement, une chaîne de caractère ou 
l'écriture d'un message décrivant l'erreur.

\subsection{Manipulation des types connus par AIDA}
\begin{description}
  \item[AidaTypCode()] : sert à reconnaître un nom de format : reçoit une
    chaîne de caractère, et renvoie le code correspondant (membre de l'enum
    \texttt{AIDA\_TYPE\_CODE)}.
  \item[AidaTypNom()] : fournit les noms (chaîne de caract. ) des types
    standard connus de AIDA (symétrique de AidaTypCode()).
  \item[AidaTypSize()] : fournit la taille (en octet) des types standard
    connus de AIDA.
\end{description}

\subsection{Manipulation des "AcqId" (noms uniques de chaque série de mesure)}
\begin{description}
  \item[Full2Id() :] construit un "Identificateur d'acquisition" (AcqId)
à partir des informations unitaires (nom de plate-forme, année, numéro
d'acquisition dans l'année).
  \item[Id2DataPath() :] construit le chemin du répertoire contenant les
données d'une acquisition.
  \item[Id2Full() :] fournit les informations unitaires (nom de plate-forme,
 année, numéro d'acquisition dans l'année) à partir de l'identificateur
d'acquisition" (AcqId).
\end{description}

\subsection{Utilisation des variables d'environnement}
         Le principe général est qu'une seule variable 
d'environnement (\texttt{AIDA\_HOME}) est nécessaire, mais
que d'autres variables (\texttt{AIDA\_LIB},  \texttt{AIDA\_DATA}, ...)
permettent de travailler sur des répertoires particuliers. La fonction
\texttt{AidaGetEnv()} gère automatiquement ce système
de valeurs particulières qui peuvent remplacer une valeur
par défaut.

\subsection{Travail sur les en-têtes}
            fonctions GetHdrValxxx, SetHdrValxxx

\subsection{Travail sur les données}
            fonctions AccessVectData, AllocVectData, FreeVectData

\section{Écriture d'un opérateur}
\label{sec:ecr_oper}

\textbf{\textit{note : le développement d'opérateurs en langage C est
abandonné au profit du développement d'opérateurs en langage Python (voir :
\nameref{sec:adc_oper}, page \pageref{sec:adc_oper}) \\
Cependant les opérateurs développés en langage C sont exécutables par
l'interpréteur de commande \texttt{adc} (voir \url{user.pdf})}}
\bigbreak

   Un opérateur, dans l'environnement AIDA, c'est une fonction
compilée, incluse dans une bibliothèque à chargement dynamique, qui travaille
sur une série temporelle de données (la notion de vecteur), avec les
fonctions d'interfaçage nécessaire pour être utilisée par l'interpréteur
de commandes adc 
(ou tout autre programme utilisant la même interface d'appel comme l'ancien
interpréteur fvb).

Concrètement, un opérateur est l'association de quatre choses :
\begin{itemize}
  \item une fonction, qui fait effectivement les calculs. Généralement en
    C ou en Fortran. Cette fonction peut faire appel à d'autres fonctions, ou
    à des bibliothèques.
  \item une fonction d'interface, qui permet à toutes les fonctions
    de pouvoir être appelées avec une même syntaxe d'appel (liste et types
    d'arguments).
  \item une fonction de documentation, qui remplit une structure décrivant
    l'opérateur et ses arguments.
  \item la documentation (format \LaTeX ).
\end{itemize}

\subsection{Le code-source d'un opérateur}
La convention de nommage est la suivante : pour un opérateur
dont le nom est <nom\_oper>, les sources sont dans un
fichier <nom\_oper>.c, qui contient la fonction d'interface
OP\_<nom\_oper> et la fonction de documentation ID\_<nom\_oper>.
La plupart du temps, la fonction qui fait le calcul est dans le
même fichier et s'appelle  <nom\_oper>, mais il peut y avoir
des exceptions lorsqu'une seule fonction peut être appelée
par deux opérateurs distincts\footnote{Par exemple, l'opérateur
Plus\_flt fait appel à la fonction de calcul Plus qui est
définie dans le source de l'opérateur Plus}
\footnote{Une autre exception est le fichier Tair.c, qui 
contient les fonction Tair, OP\_Tair et ID\_Tair, 
ainsi que OP\_Tair\_avion et ID\_Tair\_avion, qui définissent
l'opérateur Tair\_avion, mais ne font rien d'autre qu'un appel
à Tair avec une liste d'arguments spécifique}.


\subsubsection{Spécifications des propriété des vecteurs produits}
    La fonction ID\_<nom\_oper> retourne un tableau de structures,
dont chaque élément décrit un argument, et en spécifie les
attributs.

\begin{tabular}{l|l|l}
 nom          & type   & description  \\
\hline
nom	& Chaîne & \\
\hline
nom anglais	& Chaîne & \\
\hline
array\_dim	& float\_64 & \\
\hline
sample	& float\_64 & \\
\hline
NbVal	& float\_64 & \\
\hline
unit	& Chaîne & \\
\hline
ref\_time	& float\_64 & \\
\hline
misval	& float\_64 & \\
\hline
famille	& Chaîne & \\
\hline
\end{tabular}

\subsection{Écriture de la documentation d'un opérateur}

  La documentation d'un opérateur est écrite en \LaTeX, et répartie en
trois fichiers : les sections spécifiques à la description des arguments
et à la formule (ou l'algorithme) employée sont écrites dans des fichiers
séparés,
qui seront inclus dans le document général, mais qui peuvent aussi être
réutilisés et inclus dans
d'autres documents (rapports de traitement, ...)

   Comme tous les sources LaTeX, ces fichiers sont de simples fichiers
textes, modifiables avec des éditeurs généralistes (vi, emacs) ou
spécialisés (TexMaker, ...). Les
accents des textes contenus dans ces fichiers doivent être codée en
ISO-8859-15
\footnote{ISO-8859-15 est aussi baptisé
"latin9", et quasi-équivalent, pour notre usage, au
« code page » cp1252 de Windows.}.
Si vous êtes dans un environnement de travail où il est
impossible de configurer
un éditeur pour travailler en ISO-8859-15, il est toujours possible de
saisir les accents
« à l'ancienne », c'est à dire \clavier{\textbackslash'e} pour é, 
\clavier{\textbackslash`a} pour à, ...

Ces fichiers résidents dans le répertoire \pathfmt{AIDA_HOME/src/doc/oper}. Par
convention, pour un opérateur baptisé
<oper>, le fichier principal sera \pathfmt{<oper>.tex}, le fichier de
description des arguments sera \pathfmt{arguments/<oper>.args.tex},
et le fichier de description des formules sera
\pathfmt{formules/<oper>.frm.tex}. 
Enfin, pour pouvoir produire
des documentations de traitement en anglais, 
il est nécessaire qu'une version anglaise de la formule 
soit rédigée dans le fichier :
\pathfmt{formules/<oper>.frm.en.tex}.

 Il n'est jamais nécessaire de créer ces fichiers à partir de rien. Si
l'opérateur que l'on vient de créer ressemble à un opérateur existant,
il suffira de copier et modifier les fichiers de
cet opérateur.
Sinon, il existe un programme qui sait créer un squelette de
documentation : il se trouve
dans le répertoire \path{$AIDA_HOME/src/lib/oper/} et il s'appelle
\texttt{oper2doc}.

Dans le premier cas, l'effort de modification de la description peut
être minime, mais il
faut être attentif aux changements techniques (nombre, ordre et unité
des arguments).
Dans le second cas, la liste des arguments (et leurs propriétés : type,
catégorie, unité) créée
automatiquement sera parfaitement exacte, mais sans aucune intelligence
: les noms
seront Arg1, Arg2, Argn, et la formule sera vide.

La plupart du temps, il n'est pas nécessaire d'être expert en
\LaTeX{} pour écrire la documentation. Si besoin est, une
rapide introduction à l'écriture de mathématiques en
\LaTeX{}  se trouve sur le site des tuteurs de l'ENS :
\url{http://www.tuteurs.ens.fr/logiciels/latex/maths.html}.
Il faut être attentif à quelques points :\begin{itemize}
\item Les noms des variables utilisés dans la formule
doivent correspondre aux noms donnés dans la liste des
arguments.
\item Il ne faut pas utiliser la commande 
\clavier{\textbackslash\textbackslash} pour forcer 
un changement de ligne de texte, car celle-ci devient une fin 
de ligne de tableau si le fichier est inclus dans une
cellule d'un tableau. Les méthodes recommandées
sont d'utiliser le mode math hors-texte ( \verb+\[ ... \]+ ),
ou le changement de paragraphe
canonique (ligne vide).
Au pire, utiliser \clavier{\textbackslash{}newline}
\end{itemize}

\noindent Il faut ensuite :\begin{itemize}
 \item Modifier le fichier \pathfmt{CATALOG_OPER.tex} pour ajouter 
une ligne \newline
\clavier{\textbackslash input\{<nom\_oper>\}\textbackslash newpage}
dans le chapitre adapté.
 \item Produire la nouvelle
version de \pathfmt{CATALOG_OPER.pdf} par la commande \newline
\clavier{make  CATALOG\_OPER.pdf}
 \item Itérer les actions d'édition/compilation/relecture jusqu'à
  satisfaction
 \item Éventuellement, tester si la documentation qui vient d'être
écrite est bien utilisable pour générer la documentation automatique
d'un traitement. Pour cela, il faut qu'un vecteur ait été calculé
avec l'opérateur, créer un petit fichier texte contenant 
le nom de ce vecteur, puis exécuter :\newline
\clavier{mk\_doctrt <AcqId> <nom\_du\_fichier\_texte> /tmp/TEST\_DOC}
 \item Enregistrer les fichiers dans le système de gestion 
de code ("\clavier{commit}").
\end{itemize}
